%% ---------------------------------------------------------------------
%% Appendix A -- The Twelve-Week Study Plan
%% Turns the training plan of docs/core-idea.txt into a study guide that
%% points into the book.  Cross-references to chapters use ch:chNN only.
%% ---------------------------------------------------------------------
\chapter{The Twelve-Week Study Plan}
\label{ch:appA}
\index{study plan|(}

This appendix is the book's answer to a practical question: \emph{in which order, and at
what pace, should I work through these chapters if I want a working hybrid swarm planner in
three months?} It reproduces the twelve-week training plan that the book was built from,
maps every week to the chapters that teach it, and adds study notes, a capstone
description, a completion checklist and a traceability table. Read it once before you start
\cref{ch:ch01}, and return to it at the beginning of every week.

% ---------------------------------------------------------------------
\section{How to use this plan}
\label{sec:appA-intro}

\paragraph{Goal.}
The goal of the plan is to build enough theory and implementation skill to design a
\textbf{hybrid swarm planner}\index{hybrid architecture!capstone} that handles two very
different kinds of trouble: \emph{planned} conflicts between the drones you control, and
\emph{unexpected} moving obstacles such as an unknown drone crossing your formation. The
target outcome after twelve weeks is a working research prototype that combines a global
multi-agent path finding (MAPF) solver (conflict-based search, \cbs, or its
bounded-suboptimal variant \ecbs, \cref{ch:ch09,ch:ch10}), a reactive avoidance layer
(optimal reciprocal collision avoidance, \orca, or the dynamic window approach, DWA,
\cref{ch:ch13,ch:ch14}), a prediction layer (a Kalman filter or a long short-term memory
network, LSTM, \cref{ch:ch18,ch:ch20}) and a replanning layer (\dstarlite or \astar,
\cref{ch:ch04,ch:ch05}). \Cref{ch:ch24} describes how these pieces fit together and
\cref{ch:ch25} how to evaluate the result.

\begin{keyidea}[The core idea of the whole plan]
Use a global planner for planned routes, a local collision-avoidance method for surprises,
and a replanner only when the local response is insufficient. Every week of the plan builds
one of these three abilities, or the tracking and prediction that feeds them.
\end{keyidea}

\paragraph{Study load.}
Plan for approximately 6--8 hours per week. The plan prioritizes coding over passive
reading, and so does this appendix: a good split is at most two hours of reading and four
to six hours at the keyboard. Each week has a \emph{coding exercise} and a \emph{milestone}.
The milestone is a sentence you should be able to say truthfully on the last day of the
week; if you cannot, spend the first hour of the next week closing the gap rather than
moving on.
\index{study load}

\paragraph{A weekly rhythm.}
On the first day, read the week's chapters up to the worked example and trace that example
by hand. On the next three days, write the coding exercise, using the chapter's code file as
an oracle. On the fifth day, compare your implementation with the reference on the same
instance, write the milestone sentence into a lab notebook, and note what you would do
differently; keep the weekend as a buffer. The notebook matters more than it looks: the
capstone of Week~12, and any paper you write afterwards, will draw on the numbers, failure
cases and plots you produce in Weeks~1--11.
\index{lab notebook}

\paragraph{Suggested stack.}
The plan assumes the following tools. All of them are free, and the book's own code uses
only the first two.
\begin{itemize}
  \item \textbf{Python 3.10+ and NumPy.} Every reference implementation in this book is
    plain Python with NumPy (SciPy only for the QP and MILP solvers of
    \cref{ch:ch21,ch:ch22}). Write your own code the same way; you will read and debug it
    faster.
  \item \textbf{NetworkX} \cite{hagberg2008networkx} for quick graph experiments and for
    checking your graph code against an independent shortest-path implementation.
  \item \textbf{Matplotlib} for the pictures the exercises ask for: \Open and \Closed sets,
    trajectories, velocity obstacles, benchmark curves. A plot you look at every ten
    minutes is worth more than a printout you look at once.
  \item \textbf{PyBullet or gym-pybullet-drones} \cite{panerati2021learning} for a 3D
    simulator with quadrotor dynamics. You need it from Week~8 onwards and for the
    capstone; until then a 2D NumPy simulator is enough and much faster to iterate on.
  \item \textbf{Your own MAPF codebase}, if you have one. The plan was written for a reader
    who already owns a swarm or MAPF simulator; reuse its grid, instance and visualization
    code so that each week's exercise is an addition, not a rewrite.
\end{itemize}
\index{Python stack}

\begin{notebox}{How the book's code files support the exercises}
Every chapter ships a Python file \code{code/chNN\_*.py} (\code{NN} is the chapter number),
and every listing in a chapter is a verbatim excerpt of that file. Each file ends with a
self-test under \code{if \_\_name\_\_ == "\_\_main\_\_":} that runs in a few seconds and
fails loudly with an \code{assert} if the implementation is wrong. Use the files in this
order: \emph{run} the self-test to see that the reference works; \emph{write} your own
version for the week's exercise without looking at the reference; then \emph{compare} both
implementations on the same instance. The reference is your oracle, not your starting
point. When the two disagree, the chapter's worked example, whose numbers were produced by
that same file, tells you which one is wrong.
\index{code files!self-test}
\end{notebox}

% ---------------------------------------------------------------------
\section{The schedule at a glance}
\label{sec:appA-schedule}
\index{study plan!schedule}

\Cref{tab:appA-schedule} is the twelve-week schedule of the training plan, with a column
added that names the chapters to read. The \emph{Learn}, \emph{Coding exercise} and
\emph{Milestone} columns are the plan's own words. The chapters in the \emph{Read} column
are listed in the order in which you should read them; a chapter in parentheses is
background or optional depth.

{\footnotesize
\setlength{\tabcolsep}{4pt}
\begin{longtable}{@{}L{0.55cm}L{2.05cm}L{1.9cm}L{2.95cm}L{3.4cm}L{3.3cm}@{}}
\caption{The twelve-week schedule, mapped to the chapters of this book.}
\label{tab:appA-schedule}\\
\toprule
\textbf{Wk} & \textbf{Topic} & \textbf{Read} & \textbf{Learn} & \textbf{Coding exercise} & \textbf{Milestone}\\
\midrule
\endfirsthead
\multicolumn{6}{@{}l}{\textit{\Cref{tab:appA-schedule} (continued)}}\\
\toprule
\textbf{Wk} & \textbf{Topic} & \textbf{Read} & \textbf{Learn} & \textbf{Coding exercise} & \textbf{Milestone}\\
\midrule
\endhead
\midrule
\multicolumn{6}{r@{}}{\textit{continued on the next page}}\\
\endfoot
\bottomrule
\endlastfoot
1 & Foundations: Dijkstra and \astar
  & \cref{ch:ch03,ch:ch04} (\cref{ch:ch01,ch:ch02})
  & Graph search, admissible/consistent heuristics, space-time representation.
  & Implement grid \astar. Add time as a state variable. Visualize open/closed sets.
  & You can explain exactly why \astar returns an optimal path under the required conditions.\\
\addlinespace
2 & Dynamic replanning: \lpastar and \dstarlite
  & \cref{ch:ch05} (\cref{ch:ch06})
  & Incremental search and what changes when obstacles appear after planning.
  & Implement or adapt \dstarlite. Change obstacles during execution and compare runtime with rerunning \astar.
  & Your agent can repair a path rather than restarting every search.\\
\addlinespace
3 & MAPF basics and conflict types
  & \cref{ch:ch07,ch:ch08}
  & Vertex conflicts, edge/swap conflicts, time-indexed paths, sum-of-costs and makespan.
  & Implement prioritized planning for 2--5 agents. Create deliberate conflict cases.
  & You can represent and detect MAPF conflicts correctly.\\
\addlinespace
4 & Conflict-Based Search (\cbs)
  & \cref{ch:ch09}
  & Constraint tree, high-level search, low-level \astar, branching on conflicts.
  & Implement basic \cbs from scratch for a grid. Test 2, 4, and 8 agents.
  & A working optimal \cbs solver on small instances.\\
\addlinespace
5 & \ecbs and practical MAPF
  & \cref{ch:ch10} (\cref{ch:ch11})
  & Bounded suboptimality, focal search, performance/optimality trade-off.
  & Add \ecbs or reproduce an open implementation. Benchmark against \cbs.
  & You can measure the cost of faster suboptimal solutions.\\
\addlinespace
6 & Velocity Obstacles, RVO, and \orca
  & \cref{ch:ch12,ch:ch13}
  & Collision cones, relative velocity, reciprocal avoidance.
  & Build a 2D simulation with agents crossing. Compare no avoidance, VO, and \orca.
  & You understand how a local velocity choice prevents a predicted collision.\\
\addlinespace
7 & DWA and Artificial Potential Fields
  & \cref{ch:ch14,ch:ch15}
  & Short-horizon control, local minima, oscillation, robot dynamics.
  & Implement DWA or APF in a continuous environment. Create failure cases.
  & You know when local methods work and when they fail.\\
\addlinespace
8 & \rrt, \rrtstar, and continuous 3D planning
  & \cref{ch:ch16,ch:ch17}
  & Sampling, steering, collision checks, path quality in continuous space.
  & Implement \rrtstar or use a library in a 3D obstacle field. Compare to grid \astar.
  & You can choose between graph search and sampling-based planning.\\
\addlinespace
9 & Tracking unknown moving obstacles
  & \cref{ch:ch18,ch:ch19}
  & Kalman/EKF tracking, state estimation, uncertainty.
  & Simulate an external drone observed with noisy position/velocity measurements. Track it with a Kalman filter.
  & Your planner receives a stable estimated state instead of raw noisy observations.\\
\addlinespace
10 & Trajectory prediction
  & \cref{ch:ch20}
  & Constant-velocity baselines, LSTM prediction, horizon and uncertainty.
  & Compare constant velocity vs LSTM on trajectory sequences. Measure ADE/FDE or comparable prediction errors.
  & You can quantify whether learning actually improves prediction.\\
\addlinespace
11 & MPC and formation/communication constraints
  & \cref{ch:ch21,ch:ch23} (\cref{ch:ch22})
  & Receding-horizon optimization, hard/soft constraints, formation error, communication radius.
  & Create an MPC or simplified constrained controller that follows a path while preserving a formation or distance limit.
  & Constraints are explicitly enforced rather than handled only after violation.\\
\addlinespace
12 & Hybrid system integration
  & \cref{ch:ch24,ch:ch25}
  & Global MAPF + reactive avoidance + prediction + replanning.
  & Build the capstone architecture (\cref{sec:appA-capstone}) and run controlled experiments.
  & A research prototype and an experiment matrix suitable for a paper/thesis chapter.\\
\end{longtable}
}

% ---------------------------------------------------------------------
\section{The twelve weeks on one timeline}
\label{sec:appA-timeline}

The plan groups its algorithms into six families, labeled A to F, and asks for a different
depth in each: \emph{learn deeply} (A, B, C), \emph{know well} (D, E) and \emph{learn enough
to integrate} (F). \Cref{fig:appA-timeline} shows how the twelve weeks are distributed over
the families. Notice that eight of the twelve weeks go to the three ``learn deeply''
families, and that the capstone gets a full week of its own. If you are short of time, the
places to compress are Week~7 (do DWA \emph{or} APF, not both), Week~8 (use a library for
\rrtstar) and the MILP part of Week~11 (read \cref{ch:ch22}, do not implement it). Do not
compress Weeks~1, 4, 6, 9 or 12: each of them builds a layer of the capstone that the later
weeks assume.
\index{study plan!timeline}

\begin{figure}[htb]
\centering
\inputfigure{appA/timeline}
\caption{The twelve weeks grouped by algorithm family. The length of a bar is the number of
weeks the plan gives to the family; the label names the chapters read in those weeks
(\cref{ch:ch01,ch:ch02} are the foundations read in Week~1). Families A--C, which the plan
asks you to learn deeply, take eight of the twelve weeks; the capstone has a week of its
own.}
\label{fig:appA-timeline}
\end{figure}

% ---------------------------------------------------------------------
\section{Study notes, week by week}
\label{sec:appA-notes}

The notes below add to each row of \cref{tab:appA-schedule} what a table cell cannot say:
where to put your attention, what to build first, how to know that you are done, which
mistakes cost the most time, and which chapter exercises to attempt. The \emph{load}
rating, from \difficulty{1} (light) to \difficulty{3} (heavy), estimates how much of the
6--8 hours the coding exercise takes for a reader who meets the prerequisites of
\cref{ch:ch02}; plan lighter weeks around heavier ones. ``The coding exercise of'' a chapter
means the exercise marked \emph{Coding} at the end of that chapter, which is the plan's
exercise for that week.

\begin{pitfall}[Moving on with a red test]
The most expensive mistake in a self-study plan is to carry an unfinished week forward. A
space-time \astar that mishandles waiting will break the prioritized planner of Week~3,
the \cbs low level of Week~4 and the replanning layer of Week~12, and each time it will
look like a new bug. If the milestone sentence is not true on the last day of the week,
spend the first hours of the next week on it and shorten that week's exercise instead;
\cref{sec:appA-timeline} names the parts that can be compressed.
\end{pitfall}

\subsection{Week 1: Foundations -- Dijkstra and \astar}
\label{sec:appA-week01}
\index{study plan!week 01@week 1}
\noindent\textbf{Load}\ \difficulty{2}\quad\textbf{Read}\ \cref{ch:ch03,ch:ch04}, with
\cref{ch:ch01,ch:ch02} for the vocabulary and the toolbox.

\textbf{Focus.} Dijkstra's settled-node invariant (a node popped from \Open with the
smallest $\gcost$ has its final cost), and what the heuristic buys in \astar: an admissible
one makes the first expansion of the goal optimal, a consistent one makes a \Closed set safe.
Give real time to the space-time state $(v,t)$ with a wait action; it is the low-level search
of every MAPF solver in Part~III.

\textbf{Build.} Grid \astar with a pluggable heuristic (Manhattan, octile, Euclidean, zero);
the same search over $(v,t)$ with a wait action and vertex and edge constraints; a plot that
colors cells by their \Open/\Closed status.

\textbf{Done when.} You can state the optimality theorem with its conditions and say what
breaks when each is dropped; your space-time search waits or detours around a cell reserved
at a given time; with the zero heuristic both searches expand the same nodes as Dijkstra.

\textbf{Traps.} Unit diagonal cost on an 8-connected grid makes the octile and Euclidean
heuristics inadmissible; a space-time search without a time bound never terminates when the
goal is unreachable; a search that stops the instant it reaches the goal cannot honor a
later constraint on the goal cell.

\textbf{Exercises.} The coding exercise of \cref{ch:ch04} is the deliverable. Add the coding
exercise of \cref{ch:ch03}; that chapter's backward Dijkstra becomes the low-level heuristic
of \cbs in Week~4.

\subsection{Week 2: Dynamic replanning -- \lpastar and \dstarlite}
\label{sec:appA-week02}
\index{study plan!week 02@week 2}
\noindent\textbf{Load}\ \difficulty{3}\quad\textbf{Read}\ \cref{ch:ch05}; \cref{ch:ch06}
if time permits.

\textbf{Focus.} Each vertex carries $\gcost$ and $\rhs$ and is locally inconsistent when they
differ; keys, UpdateVertex and ComputeShortestPath only repair the inconsistent vertices, in
the right order. \dstarlite searches backward from the goal so that a moving start needs
only the key modifier $k_m$ instead of a re-sort of the queue.

\textbf{Build.} Turn the Week-1 search into \dstarlite; move the agent one step at a time
while obstacles appear and disappear ahead of it; count expansions and time per repair
against a fresh \astar from the current position.

\textbf{Done when.} Expansions after a small change are a small fraction of a fresh \astar,
and the path cost after every repair equals the fresh \astar cost; that equality is the
correctness test.

\textbf{Traps.} Not adding to $k_m$ when the start moves, which leaves stale keys in the
queue; updating only one endpoint of a changed edge; timing grids so small that \python
overhead dominates (count expansions too); changing the map inside ComputeShortestPath.

\textbf{Exercises.} The coding exercise of \cref{ch:ch05}; then, if time permits, the coding
exercise of \cref{ch:ch06}, since an anytime search gives the replanning layer a path within
a time budget, which \dstarlite alone does not promise.

\subsection{Week 3: MAPF basics and conflict types}
\label{sec:appA-week03}
\index{study plan!week 03@week 3}
\noindent\textbf{Load}\ \difficulty{2}\quad\textbf{Read}\ \cref{ch:ch07,ch:ch08}.

\textbf{Focus.} The formal problem of \cref{ch:ch07}: agents, time steps, wait actions, the
conflict types (vertex, edge or swap, following, cycle) and the two objectives (sum of costs,
makespan); a swap is invisible if you compare positions only. Then \cref{ch:ch08}:
prioritized planning as one space-time \astar per agent against a reservation table, fast
but incomplete and suboptimal in general.

\textbf{Build.} A conflict checker that lists every conflict of a set of time-indexed paths
with its type and time; prioritized planning for 2--5 agents; three deliberate conflict
instances (head-on in a corridor, a crossing, an agent parked on another agent's route).

\textbf{Done when.} The checker finds every planted conflict and nothing else; the planner
solves the well-formed instances; you own an instance that one priority order fails and
another solves.

\textbf{Traps.} Off-by-one in time (the position at time~0 is the start; a swap compares
$(u,v)$ at $t$ with $(v,u)$ at $t$); an agent that has arrived stays at its goal and must be
reserved there for all later times; a reservation table that stores vertices but not edges.

\textbf{Exercises.} The coding exercises of \cref{ch:ch07,ch:ch08}. Keep the checker: it is
the conflict detector of your \cbs high level next week and the ``re-check inter-agent
conflicts'' step of the capstone.

\subsection{Week 4: Conflict-Based Search}
\label{sec:appA-week04}
\index{study plan!week 04@week 4}
\noindent\textbf{Load}\ \difficulty{3}\quad\textbf{Read}\ \cref{ch:ch09}.

\textbf{Focus.} The constraint tree: a node holds constraints, one path per agent and the sum
of costs; the high level expands the cheapest node, takes its first conflict and creates two
children, each forbidding it for one agent; the low level replans only that agent with the
space-time \astar of Week~1. Optimality follows because every conflict-free solution
satisfies the constraint of at least one child.

\textbf{Build.} \cbs from scratch on a grid: vertex constraints (agent, vertex, time), edge
constraints (agent, edge, time), a constrained low level and the high-level search; test
with 2, 4 and 8 agents on an open grid and on the Week-3 corridors.

\textbf{Done when.} Your plans pass the Week-3 checker; the sum of costs equals that of the
reference implementation of \cref{ch:ch09} on every instance; you can draw the constraint
tree of a two-agent instance by hand and match it against your program's log.

\textbf{Traps.} An edge constraint has a direction and a time; a low level that returns when
it first pops the goal ignores later constraints on the goal cell; sorting the high level by
makespan while reporting sum of costs; missing conflicts after the shorter path has ended.

\textbf{Exercises.} The coding exercise of \cref{ch:ch09}. Read about cardinal conflicts and
bypass in the same chapter, but do not implement them yet; they matter when eight agents
become twenty.

\subsection{Week 5: \ecbs and practical MAPF}
\label{sec:appA-week05}
\index{study plan!week 05@week 5}
\noindent\textbf{Load}\ \difficulty{2}\quad\textbf{Read}\ \cref{ch:ch10}; \cref{ch:ch11}
for contrast.

\textbf{Focus.} Bounded suboptimality: a focal search may expand any node with
$\fcost \le w \cdot \fcost_{\min}$ and picks the one with the fewest conflicts; \ecbs does
this at both levels and guarantees a cost within a factor $w$ of optimal. Know where the
lower bound comes from (each agent's low-level $\fcost_{\min}$, summed).

\textbf{Build.} Turn your \cbs into \ecbs, or reproduce an open implementation as the plan
allows; benchmark both for $w \in \set{1.0, 1.1, 1.5, 2.0}$: success rate under a time
limit, runtime, high-level nodes and sum of costs, over several seeds.

\textbf{Done when.} A plot of cost against runtime as $w$ varies, with $w = 1$ reproducing
\cbs exactly, and a paragraph that says with numbers what suboptimality buys.

\textbf{Traps.} Using the current node's $\fcost$ instead of the minimum over \Open as the
focal bound; taking the returned path cost instead of $\fcost_{\min}$ as the lower bound;
not growing the focal list when $\fcost_{\min}$ rises; concluding from one instance.

\textbf{Exercises.} The coding exercise of \cref{ch:ch10}. Read \cref{ch:ch11} to know how
\mstar and Push-and-Swap differ from \cbs; the plan asks for awareness, not an
implementation.

\subsection{Week 6: Velocity Obstacles, RVO and \orca}
\label{sec:appA-week06}
\index{study plan!week 06@week 6}
\noindent\textbf{Load}\ \difficulty{3}\quad\textbf{Read}\ \cref{ch:ch12,ch:ch13}.

\textbf{Focus.} Everything is relative: relative position, relative velocity and a Minkowski
disc whose radius is the sum of the two radii. The velocity obstacle (VO) is the set of
relative velocities that reach the disc within the horizon $\ttc$; the reciprocal velocity
obstacle (RVO) shares the avoidance between cooperating agents; \orca turns it into one
half-plane per neighbor, so that the new velocity is a small linear program.

\textbf{Build.} A 2D simulation of crossing agents (two head-on, then eight on a circle
exchanging places) in three modes, no avoidance, VO and \orca, recording collisions,
minimum separation and time to goal.

\textbf{Done when.} No avoidance collides, VO avoids but oscillates or detours widely, \orca
is smooth and collision-free at the same speeds, and you can draw the cone and the \orca
half-plane for one pair.

\textbf{Traps.} The orientation of the half-plane (it must exclude the velocity obstacle);
the truncated cone for finite $\ttc$ is a disc plus two tangent legs, not a triangle; when
the linear program is infeasible in a crowd, choose the least-violating velocity rather than
stopping; symmetric instances deadlock unless you break the symmetry.

\textbf{Exercises.} The coding exercises of \cref{ch:ch12,ch:ch13}. Keep the simulator;
Week~7 and the capstone reuse it.

\subsection{Week 7: DWA and Artificial Potential Fields}
\label{sec:appA-week07}
\index{study plan!week 07@week 7}
\noindent\textbf{Load}\ \difficulty{2}\quad\textbf{Read}\ \cref{ch:ch14,ch:ch15}.

\textbf{Focus.} DWA searches velocity space: the dynamic window holds the velocities
reachable within one control interval, admissible velocities are those from which the robot
can still stop, and a weighted sum of heading, clearance and velocity terms picks the
command. Artificial potential fields (APF) have no search, only forces, attractive to the
goal and repulsive from obstacles, which is why APF is fast and why it has local minima and
oscillation.

\textbf{Build.} One of the two (DWA if your drones have real acceleration limits, APF as a
two-hour baseline) in the Week-6 simulator; then failure cases: a U-shaped obstacle, a
narrow passage, a dead end that a short DWA horizon cannot see.

\textbf{Done when.} A picture of each failure with the reason, and a tuning table of how the
DWA weights or the APF gains change the behavior.

\textbf{Traps.} DWA terms of different units summed without normalization; skipping the
braking (admissibility) check; repulsive forces that grow without bound near obstacles
(clamp them); comparing the methods at different control rates.

\textbf{Exercises.} The coding exercise of \cref{ch:ch14} or of \cref{ch:ch15}, plus the
conceptual exercises of the other chapter. DWA is the plan's alternative to \orca in the
local safety layer.

\subsection{Week 8: \rrt, \rrtstar and continuous 3D planning}
\label{sec:appA-week08}
\index{study plan!week 08@week 8}
\noindent\textbf{Load}\ \difficulty{2}\quad\textbf{Read}\ \cref{ch:ch16,ch:ch17}.

\textbf{Focus.} Sample, nearest, steer and collision-free plus a goal bias make \rrt; the
near set, parent choice and rewiring with a shrinking connection radius make \rrtstar
asymptotically optimal; informed sampling restricts samples to the ellipsoid that can still
improve the solution. Know what probabilistic completeness and asymptotic optimality
promise, and that neither says anything about finite time.

\textbf{Build.} \rrtstar, yours or a library's, in a 3D field of box or sphere obstacles;
grid \astar on the same field at two or three resolutions; compare path length, time and
memory.

\textbf{Done when.} A convergence plot of the \rrtstar cost against iterations, a table of
\astar per resolution against \rrtstar per sample budget, and a one-paragraph rule for when
you would choose each.

\textbf{Traps.} Checking the new vertex but not the segment to it; a quadratic
nearest-neighbor search (use a KD-tree above a few thousand nodes); a rewire that changes a
parent without updating the descendants' costs; comparing raw paths without smoothing
either.

\textbf{Exercises.} The coding exercises of \cref{ch:ch16,ch:ch17}. The 3D field you build
here is the map of the capstone.

\subsection{Week 9: Tracking unknown moving obstacles}
\label{sec:appA-week09}
\index{study plan!week 09@week 9}
\noindent\textbf{Load}\ \difficulty{2}\quad\textbf{Read}\ \cref{ch:ch18,ch:ch19}.

\textbf{Focus.} Predict with the motion model (the state moves, the covariance grows),
update with the measurement (the gain weighs model against measurement). The
constant-velocity model with white-noise-acceleration process noise is the workhorse; gating
rejects outliers; a missing measurement is a predict without update. \Cref{ch:ch19} adds
the extended and unscented Kalman filters (EKF, UKF) and the particle filter for targets
that turn.

\textbf{Build.} An external drone on a curved path, observed with noisy position (optionally
velocity) measurements at 10--20~Hz with dropouts; a Kalman filter with tuned $\mat{Q}$ and
$\mat{R}$; the error of the estimate against the truth next to the error of the raw
measurements.

\textbf{Done when.} The estimate beats the raw measurements, the velocity estimate
extrapolates one or two seconds ahead without jitter, and the Week-6 simulator consumes the
estimate and its covariance instead of the observation.

\textbf{Traps.} $\mat{R}$ below the real noise (the filter trusts noise) or $\mat{Q}$ too
small (it lags on turns and becomes overconfident); a wrong $\dt$ in the transition matrix;
trusting the covariance when the model is violated (check the normalized innovations);
tuning on one trajectory.

\textbf{Exercises.} The coding exercise of \cref{ch:ch18}; the coding exercise of
\cref{ch:ch19} if your intruder turns sharply.

\subsection{Week 10: Trajectory prediction}
\label{sec:appA-week10}
\index{study plan!week 10@week 10}
\noindent\textbf{Load}\ \difficulty{3}\quad\textbf{Read}\ \cref{ch:ch20}.

\textbf{Focus.} Average and final displacement error (ADE, FDE) as the yardsticks;
constant-velocity and Kalman baselines; the LSTM as a sequence model that reads past
positions and writes future ones; horizon, uncertainty and the evaluation protocol.

\textbf{Build.} Trajectories from the Week-6 simulator (agents that avoid one another are
not constant-velocity, which is what makes learning worthwhile), split by trajectory rather
than by time step; both baselines; an LSTM (a deep-learning framework is fine; the book's
own code stays in NumPy); ADE/FDE at 1, 2 and 3~s.

\textbf{Done when.} A table of ADE/FDE per horizon and method over several training seeds,
and an honest conclusion: ``the LSTM did not beat constant velocity on this data'' is a
legitimate and common result at short horizons.

\textbf{Traps.} Leakage between training and test trajectories; predicting absolute
coordinates instead of displacements in the agent's frame; a weak baseline that uses two raw
points instead of a smoothed velocity; comparing at different horizons; tuning on the test
set.

\textbf{Exercises.} The coding exercise of \cref{ch:ch20}. Save the model and the
baselines; the ``prediction quality'' axis of the capstone uses all of them.

\subsection{Week 11: MPC and formation/communication constraints}
\label{sec:appA-week11}
\index{study plan!week 11@week 11}
\noindent\textbf{Load}\ \difficulty{3}\quad\textbf{Read}\ \cref{ch:ch21,ch:ch23};
\cref{ch:ch22} as background.

\textbf{Focus.} Receding horizon: solve a short optimal control problem, apply the first
input, repeat. With the double-integrator model and a quadratic cost it is a quadratic
program; velocity and acceleration bounds are linear; separation constraints are nonconvex
and are linearized around the previous solution; slack variables keep the program feasible.
From \cref{ch:ch23}, displacement-based formation and the formation error. Mixed-integer
linear programming (MILP, \cref{ch:ch22}) is the exact alternative; read it for the cost of
exactness.

\textbf{Build.} A model predictive controller (MPC), or the simplified constrained
controller the plan allows, that follows a nominal path while keeping the formation error
bounded or every pairwise distance below the communication radius; log violations and slack.

\textbf{Done when.} Zero violations with hard constraints, logged and small slack with soft
ones, a plot of the formation error over time, and a visible difference when the constraint
is switched off.

\textbf{Traps.} A horizon too long for the time budget or too short to see a conflict
coming; hard constraints that turn infeasible when an intruder closes in (use slack);
linearizing around a stale trajectory; a model $\dt$ that differs from the control interval.

\textbf{Exercises.} The coding exercises of \cref{ch:ch21,ch:ch23}; the coding exercise of
\cref{ch:ch22} is optional.

\subsection{Week 12: Hybrid system integration}
\label{sec:appA-week12}
\index{study plan!week 12@week 12}
\noindent\textbf{Load}\ \difficulty{3}\quad\textbf{Read}\ \cref{ch:ch24,ch:ch25}.

\textbf{Focus.} The architecture and decision logic of \cref{sec:appA-capstone}: the safety
horizon and the time-to-collision trigger, reconnection to the nominal path, the replanning
triggers, the conflict re-check, and the experiment matrix and metrics of \cref{ch:ch25}.

\textbf{Build.} The capstone: the four layers you already own, connected behind clean
interfaces, and an experiment runner that takes a scenario and a seed and writes one row of
metrics per run.

\textbf{Done when.} The experiment matrix runs unattended from one script, every run is
reproducible from its seed, and a table and plots compare local-only, replan-only and hybrid
strategies under the same conditions.

\textbf{Traps.} One file that mixes the layers; different time steps for the discrete global
plan and the continuous local layer (convert in one place); measuring success rate only;
running experiments before each layer passes its own tests.

\textbf{Exercises.} The coding exercises of \cref{ch:ch24,ch:ch25}.

% ---------------------------------------------------------------------
\section{The capstone: a hybrid collision-avoidance prototype}
\label{sec:appA-capstone}
\index{capstone}

The capstone is the plan's target outcome: a research prototype that combines a global
multi-agent planner, a prediction layer, a local safety layer and a replanning layer.
Everything you built in Weeks~1--11 is a component of it. \Cref{ch:ch24} designs the
architecture in full, with its state machine, its safety horizon and a discussion of which
guarantees survive; \cref{ch:ch25} shows how to evaluate it. This section restates the
plan's specification so that you can check your prototype against it.

\subsection{The four layers}
\label{sec:appA-layers}
\Cref{tab:appA-layers} lists the four layers in the plan's words. The numbering is the
plan's; it is neither the order in which the layers act at run time (that is the decision
logic below) nor the order in which you build them (that is the schedule).

\begin{table}[htb]
\centering\small
\caption{The four layers of the capstone architecture, in the plan's words, with the
chapters that teach each layer and the weeks in which you build it.}
\label{tab:appA-layers}
\begin{tabular}{@{}L{3.1cm}L{7.3cm}L{2.3cm}c@{}}
\toprule
\textbf{Layer} & \textbf{Job} & \textbf{Chapters} & \textbf{Weeks}\\
\midrule
1\ Global swarm planner (\cbs/\ecbs)
  & Plan nominal, time-indexed paths for the drones under your control. Resolve conflicts among connected swarm members before execution.
  & \cref{ch:ch09,ch:ch10} & 3--5\\
\addlinespace
2\ Prediction layer (Kalman filter / LSTM)
  & Observe an unknown external drone. Estimate its current state and predict a short future trajectory with uncertainty.
  & \cref{ch:ch18,ch:ch20} & 9--10\\
\addlinespace
3\ Local safety layer (\orca or DWA)
  & When a predicted collision enters a short safety horizon, generate a temporary safe velocity or local maneuver without immediately discarding the global route.
  & \cref{ch:ch13,ch:ch14} & 6--7\\
\addlinespace
4\ Replanning layer (\dstarlite or \astar)
  & If the local maneuver cannot safely return the drone to its nominal route, replan from the current state. Then re-check swarm conflicts and formation/communication constraints.
  & \cref{ch:ch04,ch:ch05} & 1--2, 11\\
\bottomrule
\end{tabular}
\end{table}

\subsection{The decision logic}
\label{sec:appA-logic}
The plan suggests five steps for every controlled drone, at every control cycle.
\Cref{ch:ch24} turns them into a state machine with explicit triggers (a time to collision
below a threshold, a reconnection cost above a bound, a violated constraint) and explains
what happens to the optimality of the \cbs plan once a drone leaves it.
\begin{enumerate}
  \item Follow the nominal \cbs/\ecbs path while no predicted conflict is inside the safety
    horizon.
  \item If an unknown moving object creates a near-term risk, invoke \orca/DWA to produce a
    local avoidance action.
  \item After the conflict clears, attempt to reconnect to the nominal path at a safe future
    waypoint.
  \item If reconnection is impossible, too costly, or violates formation/communication
    constraints, run \dstarlite/\astar from the current state.
  \item Re-check inter-agent conflicts after any significant replan.
\end{enumerate}
\index{decision logic!capstone}

\begin{notebox}{Keep the layers separable}
Give each layer one entry point: the global planner takes an instance and returns
time-indexed paths; the predictor takes observations and returns a predicted trajectory
with covariance; the local layer takes states and predictions and returns a velocity; the
replanner takes a start state and returns a path. The ``avoidance strategy'' axis of the
experiments is then a switch, not a rewrite: local-only disables layer~4, replan-only
disables layer~3, the hybrid uses both.
\end{notebox}

\subsection{Experiments and metrics}
\label{sec:appA-experiments}
\Cref{tab:appA-experiments} lists the factors the plan asks you to vary; \cref{ch:ch25}
gives the scenario generator and the statistics to report. As a full factorial design the
five factors have $4 \times 4 \times 4 \times 3 \times 4 = 768$ cells before seeds; you will
not run all of them. Fix every factor at a default, vary one or two at a time, and let the
research question of \cref{ch:ch24} decide which.
\index{experiment matrix}

\begin{table}[htb]
\centering\small
\caption{The experiment factors of the capstone and their levels, in the plan's words.}
\label{tab:appA-experiments}
\begin{tabular}{@{}L{3.3cm}L{7.4cm}L{3.6cm}@{}}
\toprule
\textbf{Factor} & \textbf{Levels} & \textbf{Where in this book}\\
\midrule
Controlled drones & 2, 4, 8, and more if computationally feasible & \cref{ch:ch25}; scaling of \cbs/\ecbs in \cref{ch:ch09,ch:ch10}\\
\addlinespace
Unknown external drones & 0, 1, 2, and multiple crossing trajectories & \cref{ch:ch25}\\
\addlinespace
Prediction quality & perfect state, noisy state, constant-velocity prediction, learned prediction & \cref{ch:ch18,ch:ch20,ch:ch25}\\
\addlinespace
Avoidance strategy & local-only, global-replan-only, and hybrid local + replanning & \cref{ch:ch24,ch:ch25}\\
\addlinespace
Constraints & no formation requirement vs formation-preserving; no communication radius vs bounded communication range & \cref{ch:ch21,ch:ch23}; during avoidance in \cref{ch:ch24}\\
\bottomrule
\end{tabular}
\end{table}

\paragraph{Metrics.}
Report, for every run, the metrics the plan lists: \textbf{collision rate}, \textbf{minimum
separation}, \textbf{path length}, \textbf{travel time}, \textbf{makespan}, \textbf{sum of
costs}, \textbf{replanning count}, \textbf{computation time}, \textbf{formation error} and
\textbf{communication violations}.\index{metrics!capstone} \Cref{ch:ch25} defines each of
them and the statistics to report over seeds; makespan and sum of costs are the MAPF
objectives of \cref{ch:ch07}, and the replanning count is the direct measure of the decision
logic of \cref{ch:ch24}, since it counts how often the local layer was not enough. Two of
them deserve a warning: a collision rate of zero hides near misses, which is what minimum
separation is for, and a mean computation time hides the cycle that missed its deadline, so
report the maximum as well.

% ---------------------------------------------------------------------
\section{Completion checklist}
\label{sec:appA-checklist}
\index{completion checklist}

The plan ends with ten statements. \Cref{tab:appA-checklist} repeats them and names the
chapters where each ability is taught. Tick a line only when you can do what it says without
opening the book; together, the ten lines are a fair summary of what a thesis committee or
a reviewer will expect you to know.

\begin{table}[htb]
\centering\small
\caption{The plan's completion checklist, with the chapters that teach each item.}
\label{tab:appA-checklist}
\begin{tabular}{@{}cL{8.3cm}L{5.4cm}@{}}
\toprule
& \textbf{Statement} & \textbf{Where in this book}\\
\midrule
$\square$ & I can implement and explain \astar and space-time \astar. & \cref{ch:ch04}; reservation tables in \cref{ch:ch08}\\
$\square$ & I can explain the \cbs constraint tree and implement vertex and edge constraints. & \cref{ch:ch09}; conflict types in \cref{ch:ch07}\\
$\square$ & I understand the speed/optimality trade-off in \ecbs. & \cref{ch:ch10}\\
$\square$ & I can explain VO/RVO/\orca using relative position and velocity. & \cref{ch:ch12,ch:ch13}\\
$\square$ & I can track an unknown drone from noisy observations. & \cref{ch:ch18,ch:ch19}\\
$\square$ & I can compare a simple prediction baseline with LSTM prediction. & \cref{ch:ch20}\\
$\square$ & I can trigger local avoidance based on a time-to-collision or safety-horizon rule. & time to collision in \cref{ch:ch12}; the trigger in \cref{ch:ch24}\\
$\square$ & I can reconnect to the global route or trigger \dstarlite/\astar replanning. & \cref{ch:ch05}; reconnection in \cref{ch:ch24}\\
$\square$ & I can enforce formation and/or communication-range constraints. & \cref{ch:ch21,ch:ch23}; during avoidance in \cref{ch:ch24}\\
$\square$ & I can evaluate the hybrid system with reproducible metrics and controlled scenarios. & \cref{ch:ch25}\\
\bottomrule
\end{tabular}
\end{table}

% ---------------------------------------------------------------------
\section{Traceability: from the plan to the chapters}
\label{sec:appA-traceability}
\index{traceability table}

\Cref{tab:appA-traceability} maps every algorithm of the training plan to the chapter that
teaches it and the week that uses it, with the plan's priority tag; where the plan qualified
a tag, the qualification is kept next to it. The plan's one-line reason for each algorithm
is in the algorithm map of \cref{ch:ch01}. Use the table in two directions: to check that the
book covers everything the plan asks for, and, when you open a chapter, to see how much
depth the plan expects of you there.

{\footnotesize
\setlength{\tabcolsep}{4pt}
\begin{longtable}{@{}L{5.0cm}L{4.6cm}L{2.8cm}C{0.8cm}@{}}
\caption{Every algorithm of the training plan, its priority, and where this book teaches it.}
\label{tab:appA-traceability}\\
\toprule
\textbf{Algorithm} & \textbf{Priority in the plan} & \textbf{Chapter} & \textbf{Wk}\\
\midrule
\endfirsthead
\multicolumn{4}{@{}l}{\textit{\Cref{tab:appA-traceability} (continued)}}\\
\toprule
\textbf{Algorithm} & \textbf{Priority in the plan} & \textbf{Chapter} & \textbf{Wk}\\
\midrule
\endhead
\midrule
\multicolumn{4}{r@{}}{\textit{continued on the next page}}\\
\endfoot
\bottomrule
\endlastfoot
\multicolumn{4}{@{}l}{\textbf{A. Single-agent graph search} -- learn deeply}\\
\addlinespace[2pt]
Dijkstra & \priority{High} & \cref{ch:ch03} & 1\\
\astar & \priority{Essential} & \cref{ch:ch04} & 1\\
\dstarlite & \priority{Essential} & \cref{ch:ch05} & 2\\
Lifelong Planning \astar (\lpastar) & \priority{Medium} & \cref{ch:ch05} & 2\\
Anytime Repairing \astar (\arastar) & \priority{Medium} & \cref{ch:ch06} & 2\\
\midrule
\multicolumn{4}{@{}l}{\textbf{B. Multi-agent path finding} -- learn deeply}\\
\addlinespace[2pt]
Prioritized Planning & \priority{High} & \cref{ch:ch08} & 3\\
Conflict-Based Search (\cbs) & \priority{Essential} & \cref{ch:ch09} & 4\\
Enhanced \cbs (\ecbs) & \priority{Essential} & \cref{ch:ch10} & 5\\
\mstar & \priority{Medium} & \cref{ch:ch11} & 5\\
Push-and-Swap / Push-and-Rotate & \priority{Awareness} & \cref{ch:ch11} & 5\\
\midrule
\multicolumn{4}{@{}l}{\textbf{C. Local/reactive collision avoidance} -- learn deeply}\\
\addlinespace[2pt]
Velocity Obstacles (VO) & \priority{Essential} & \cref{ch:ch12} & 6\\
Reciprocal Velocity Obstacles (RVO) & \priority{High} & \cref{ch:ch13} & 6\\
\orca & \priority{Essential} & \cref{ch:ch13} & 6\\
Dynamic Window Approach (DWA) & \priority{High} & \cref{ch:ch14} & 7\\
Artificial Potential Fields (APF) & \priority{High} & \cref{ch:ch15} & 7\\
\midrule
\multicolumn{4}{@{}l}{\textbf{D. Sampling-based motion planning} -- know well}\\
\addlinespace[2pt]
\rrt & \priority{High} & \cref{ch:ch16} & 8\\
\rrtstar & \priority{High} & \cref{ch:ch17} & 8\\
\irrtstar & \priority{Medium} & \cref{ch:ch17} & 8\\
\midrule
\multicolumn{4}{@{}l}{\textbf{E. Prediction and tracking} -- know well}\\
\addlinespace[2pt]
Kalman Filter & \priority{Essential} & \cref{ch:ch18} & 9\\
Extended Kalman Filter (EKF) & \priority{High} & \cref{ch:ch19} & 9\\
Unscented Kalman Filter (UKF) & \priority{Medium} & \cref{ch:ch19} & 9\\
Particle Filter & \priority{Medium} & \cref{ch:ch19} & 9\\
LSTM trajectory prediction & \priority{Essential} {\scriptsize for your planned research} & \cref{ch:ch20} & 10\\
Transformer trajectory prediction & \priority{Awareness to High} & \cref{ch:ch20} & 10\\
\midrule
\multicolumn{4}{@{}l}{\textbf{F. Optimization and control} -- learn enough to integrate}\\
\addlinespace[2pt]
Model Predictive Control (MPC) & \priority{Essential} & \cref{ch:ch21} & 11\\
Mixed-Integer Linear Programming (MILP) & \priority{Medium} & \cref{ch:ch22} & 11\\
Consensus / formation control & \priority{Essential} {\scriptsize for formation-preserving work} & \cref{ch:ch23} (\cref{ch:ch21}) & 11\\
\end{longtable}
}

% ---------------------------------------------------------------------
\section{Reading order and primary sources}
\label{sec:appA-reading}
\index{reading order}

The plan also gives a reading order for the literature, reproduced in
\cref{tab:appA-reading} together with the chapter that covers each item and the primary
source to read next to it. The first six items are the plan's highest priority: they are
the layers of the capstone plus the MAPF definitions that make the conflict re-check
precise. To read a paper next to its chapter: read the chapter up to the end of its worked
example; then read the paper's algorithm section and compare its pseudocode line by line
with the chapter's; then read the paper's experiments and ask which of its claims your
implementation of that week reproduces. Item~10 is a book rather than a paper: LaValle's
\emph{Planning Algorithms} \cite{lavalle2006planning} is the reference to keep open for
anything about configuration spaces, sampling and collision checking that
\cref{ch:ch02,ch:ch16,ch:ch17} treat only as far as this book needs.

\begin{table}[htb]
\centering\small
\caption{The plan's reading order, mapped to chapters and to the primary sources.}
\label{tab:appA-reading}
\begin{tabular}{@{}rL{3.4cm}L{3.3cm}L{7.0cm}@{}}
\toprule
\textbf{\#} & \textbf{Read} & \textbf{Chapter(s)} & \textbf{Primary source}\\
\midrule
1 & \astar & \cref{ch:ch04} & Hart, Nilsson and Raphael \cite{hart1968formal}\\
2 & \dstarlite & \cref{ch:ch05} & Koenig and Likhachev \cite{koenig2002dstarlite}\\
3 & \cbs & \cref{ch:ch09} & Sharon et al.\ \cite{sharon2015cbs}\\
4 & MAPF overview & \cref{ch:ch07} & Stern et al.\ \cite{stern2019mapf}\\
5 & Velocity Obstacles & \cref{ch:ch12} & Fiorini and Shiller \cite{fiorini1998motion}\\
6 & \orca & \cref{ch:ch13} & van den Berg et al.\ \cite{vandenberg2011orca}\\
\addlinespace
7 & DWA & \cref{ch:ch14} & Fox, Burgard and Thrun \cite{fox1997dwa}\\
8 & \rrt & \cref{ch:ch16} & LaValle \cite{lavalle1998rrt}\\
9 & \rrtstar & \cref{ch:ch17} & Karaman and Frazzoli \cite{karaman2011sampling}\\
10 & Planning reference & \cref{ch:ch02,ch:ch16,ch:ch17} & LaValle, \emph{Planning Algorithms} \cite{lavalle2006planning}\\
11 & MPC & \cref{ch:ch21} & Rawlings, Mayne and Diehl \cite{rawlings2017mpc}\\
12 & Trajectory prediction & \cref{ch:ch20} & Rudenko et al.\ (survey) \cite{rudenko2020human}; Alahi et al.\ (Social LSTM) \cite{alahi2016social}\\
\bottomrule
\end{tabular}
\end{table}

\paragraph{After Week 12.}
The plan is explicit that the individual algorithms above are established, and that the
research lies in how they are combined, constrained, predicted and evaluated: unknown
non-cooperative drones, local avoidance that does not needlessly destroy the guarantees of
the global plan, formation and communication constraints during avoidance,
prediction-aware constraints that carry uncertainty, an automatic decision between local
action and global replan, computation for larger swarms, and benchmarks that compare
local-only, replan-only and hybrid approaches under the same conditions. \Cref{ch:ch24}
discusses these directions and \cref{ch:ch25} the benchmark; the prototype and the
experiment matrix you finish in Week~12 are the starting point for both.
\index{research directions}
\index{study plan|)}
